\par \textbf{Степень расширения}
\begin{itemize}
    \item Идейно: раскладываем произвольный элемент расширения по базису или используем теорему о башнях расширений (билет 41 на уд)
    \item Примеры:
\begin{enumerate}
    \item $\Comp \supset \R$. 
    \par Базисные векторы: $i, 1 \Rightarrow$ размерность 2 $\Rightarrow$ степень расширения 2
    \item $\Q(\sqrt{2}) \supset \Q$.
    \par $\Q[\sqrt{2}]=\{a\sqrt{2}+b \: | \: a, b \in \Q\}$. Тогда $\Q(\sqrt{2})=\{\frac{a\sqrt{2}+b}{c+d\sqrt{2}} \: | \: a,b,c,d \in \Q\}$. Домножая числитель и знаменатель на $c-d\sqrt{2}$ (знаменатель не окажется равным 0 в силу иррациональности $\sqrt{2}$) получим $\Q(\sqrt{2})=\{a\sqrt{2}+b \: | \: a,b \in \Q\}$, значит степень расширения равна 2
    \item $\Q(\omega) \supset \Q$.
    \par Так как $\omega^2=-1-\omega$, то $\Q(\omega)=\{\frac{a\omega+b}{c\omega+d} \: | \: a,b,c,d \in \Q\}$. Домножим числитель и знаменатель на $c+d\overline{\omega}$. Знаменатель опять не обратится в ноль и после приведения подобных слагаемых получим $\Q(\omega)=\{\alpha\omega+\beta \: | \: \alpha, \beta \in \Q\} \Rightarrow$ степень расширения равна 2
    \item $\Q(i) \supset \Q$.
    \par $x^2 + 1$ -- минимальный многочлен, поэтому $[\Q(i):\Q]=2$
    \item $\Q(\sqrt[4]{2}, i) \supset \Q$.
    \par Рассмотрим башню расширений $\Q \subset \Q(\sqrt[4]{2})\subset \Q(\sqrt[4]{2})(i)=\Q(\sqrt[4]{2}, i)$. \newline Так как $[\Q(\sqrt[4]{2}):\Q]=4$ (минимальный многочлен $x^4-2$ - неприводим по признаку Эйзенштейна), то по теореме о башне расширений (билет 41 на уд) \newline $[\Q(\sqrt[4]{2},i):\Q]= [\Q(\sqrt[4]{2},i):\Q(\sqrt[4]{2})] \cdot [\Q(\sqrt[4]{2}):\Q] = 2 \cdot 4 = 8$
\end{enumerate}
\end{itemize}


\par \textbf{Минимальный многочлен:} 
\begin{itemize}
    \item Идейно: подбираем многочлен через возведение в степень. Наверняка попадем в минимальный. Остается только доказать что он минимальный (тупой проверкой, что меньше не может быть; доказательством неприводимости; через степень расширения)
    
    \item Примеры:
    \begin{enumerate}
        \item $\sqrt{2}$ над $\Q$
        \par $x=\sqrt{2} \Rightarrow x^2-2=0$. Если это не минимальный многочлен, то $\deg m=1$, но $\sqrt{2} \not\in \Q \Rightarrow x^2-2$ - минимальный
        \item $\sqrt[7]{5}$ над $\Q$.
        \par $x=\sqrt[7]{5}\Rightarrow x^7-5=0$. Этот многочлен неприводим по признаку Эйзентштейна для $p=5$. Пусть это не минимальный многочлен, значит он должен делиться на минимальный, но он неприводим - противоречие.
        
        \item $\sqrt{2}+\sqrt{5}$ над $\Q$
        \par $x=\sqrt{2}+\sqrt{5}\Rightarrow x^2=7+2\sqrt{10}\Rightarrow x^2-7=2\sqrt{10}\Rightarrow x^4-14x^2+9=0$. Рассмотрим 2 решения \begin{itemize}
            \item Докажем, что многочлен неприводим. Для этого разложим его на множители в $\Comp$, а затем покажем, что никакие комбинации скобок не дадут нам многочленов из $\Q[x]$
            $$x^4-14x^2+9=(x-(\sqrt{2}+\sqrt{5}))(x-(\sqrt{2}-\sqrt{5}))(x-(-\sqrt{2}+\sqrt{5}))(x-(-\sqrt{2}-\sqrt{5}))$$
            \par Каждая скобка очевидно многочлен с иррациональными коэффициентами (если $\sqrt{2}+\sqrt{5}=r\in \Q \Rightarrow 7+2\sqrt{10}=r^2 \in \Q$, но $\sqrt{10}$ - иррационален). Значит нам остается только перебрать попарные комбинации скобок. Рассмотрим комбинацию в которой содержится первая скобка. Простым перебором трех вариантов получаем, что все многочлены с иррациональными коэффициентами, а значит наш исходный многочлен неприводим
            \item Докажем, что $[\Q(\sqrt{2}+\sqrt{5}):\Q]=4$, тогда минимальный многочлен должен иметь степень 4. 
            $$(\sqrt{2}+\sqrt{5})(\sqrt{5}-\sqrt{2})=3 \;\;\Rightarrow\;\; \sqrt{5}-\sqrt{2}=\frac{3}{\sqrt{5}+\sqrt{2}} \in \Q(\sqrt{2}+\sqrt{5}) \Rightarrow$$ $$\Rightarrow 2\sqrt{5} \in \Q(\sqrt{2}+\sqrt{5}) \;\;\Rightarrow\;\; \sqrt{5} \in \Q(\sqrt{2}+\sqrt{5})$$
            \par Значит $\Q \subset \Q(\sqrt{5}) \subset \Q(\sqrt{2}+\sqrt{5})$ и по теореме о башне расширений (билет 41 на уд) выполнено
            $$[\Q(\sqrt{2}+\sqrt{5}):\Q]=[\Q(\sqrt{2}+\sqrt{5}):\Q(\sqrt{5})]\cdot [\Q(\sqrt{5}):\Q]=2\cdot 2=4$$
        \end{itemize}
    \end{enumerate}
\end{itemize}

\par \textbf{Поле разложения:}
\begin{itemize}
    \item Примеры:
    \begin{enumerate}
        \item $x^n-2=(x-\sqrt[n]{2})(x-\sqrt[n]{2}\xi_n)\ldots(x-\sqrt[n]{2}\xi_n^{n-1})$, где $\xi_n$ - примитивный корень из 1. Тогда поле разложения $\mathbb{F}=\Q(\sqrt[n]{2}, \sqrt[n]{2}\xi_n, \ldots)=\Q[\sqrt[n]{2}, \sqrt[n]{2}\xi_n, \ldots]$. Так как это поле, то можем поделить второй порождающий на первый и получить $\xi_n$. Так как остальные элементы выражаются как произведение $\xi_n$ и $\sqrt[n]{2}$, то $\mathbb{F}=\Q[\sqrt[n]{2}, \xi_n]$
    \end{enumerate}
\end{itemize} 

\par \textbf{Примитивный элемент:}
\begin{itemize}
    \item  Идейно: по теореме о примитивном элементе, любое конечное расширение $K$ поля $F$ при $char F = 0$ можно представить в виде $K=F(\gamma)$ (где $\gamma$ - примитивный элемент). При доказательстве этой теоремы было доказано утверждение, что если $K=F(\alpha, \beta)$, то $\gamma=\alpha+c\beta$, где $c\neq \frac{\alpha-\alpha_i}{\beta-\beta_i}$, где $\alpha_i \neq \alpha, \beta_i \neq \beta$ - остальные корни минимальных многочленов для $\alpha$ и $\beta$ соответственно. Таким образом, нам просто нужно найти все корни минимальных многочленов и взять $c$ подходящую под условие
    \item Примеры: 
    \begin{enumerate}
        \item $\Q(\sqrt{2})\supset \Q$. Очевидно, что ответ $\sqrt{2}$
        \item $\Q(\sqrt{2}, \sqrt{3}, \sqrt{5})\supset \Q$.
        \par Будем действовать пошагово: сначала найдем примитивный элемент для $\Q(\sqrt{2}, \sqrt{3})$. Корни минимального многочлена для $\sqrt{2}: \pm \sqrt{2}$, для $\sqrt{3}: \pm \sqrt{3}$. Значит единственное запрещенное значение $c=\frac{\sqrt{2}-(-\sqrt{2})}{\sqrt{3}-(-\sqrt{2})}=\frac{\sqrt{2}}{\sqrt{3}}$. Возьмем $c=1$. Тогда $\sqrt{2}+\sqrt{3}$ - примитивный элемент
        \par Теперь надо найти примитивный элемент для $\Q(\sqrt{2}+\sqrt{3}, \sqrt{5})$. Корни для $\sqrt{2}+\sqrt{3}: \pm \sqrt{2} \pm \sqrt{3}$, для $\sqrt{5}: \pm \sqrt{5}$. Аккуратно расписав все запрещенные значения $c$ (их 3) видим, что среди них снова нет единицы, а значит $\sqrt{2}+\sqrt{3}+\sqrt{5}$ - примитивный
    \end{enumerate}
\end{itemize}
